% An invented source paper for the Sextant demo. Nothing in it is real.
\documentclass[10pt]{article}
\usepackage[T1]{fontenc}
\usepackage[margin=2.4cm]{geometry}
\usepackage{lmodern}
\pagestyle{empty}
\begin{document}
\begin{center}
{\small\itshape Imaginary Soil Science Letters, vol.~4, pp.~33--38, 2019\par}\vspace{10pt}
{\LARGE\bfseries Straw Mulch Depth and Weed Emergence in Allotment Plots\par}\vspace{8pt}
{\large Morgan Placeholder\par}\vspace{4pt}
{\small Department of Imaginary Studies, Nowhere University\par}
\end{center}
\vspace{6pt}
\noindent\textbf{Abstract.} We compared bare soil with straw mulch laid at three depths on allotment plots and counted emerging weeds and surface moisture over one summer.

\section*{Methods}
Thirty plots were split into four treatments: bare soil, and straw laid two, five and eight centimetres deep. Weed seedlings were counted in a fixed frame every fortnight, and surface moisture was measured at midday.

\section*{Results}
Counted over the whole summer and averaged across the plots in each treatment, a straw mulch layer 8 cm deep cut weed emergence by 54\% compared with bare soil, while the 2 cm layer made no measurable difference. The deepest layer also kept the surface noticeably moister through the hottest weeks of the season.

\section*{Discussion}
Depth matters more than material: a thin layer blows away or lets light through. Straw is cheap and easy to source, which makes it a practical first step for volunteer-run plots.
\end{document}
